\section{The variational principle}

Throughout, $S = \mathbb Z^d$, $\lambda$ is a probability measure on $E$, $\Phi$ a
shift-invariant absolutely summable potential with Gibbs specification $\gamma^\Phi$, $h$ the
specific entropy of Theorem \ref{thm:sent-15-12} and $P$ the pressure of Theorem
\ref{thm:pr-15-30}.

\begin{proposition}[Georgii (15.28)--(15.30)(b), (15.33)--(15.35)]
    \label{prop:vp-15-30b}
    \uses{thm:sent-15-12, thm:pr-15-30, prop:pr-15-25, def:erg-invariant}
    \lean{Potential.relativeEntropyIn_bind_le_add, Potential.tendsto_relativeEntropyIn_bind_div_card_congr, Potential.tendsto_relativeEntropyIn_div_card, Potential.specificRelativeEntropy, Potential.tendsto_relativeEntropyIn_gibbsSpecification_div_card_zero, Potential.relativeEntropyIn_gibbsSpecification_eq, Potential.specificRelativeEntropy_nonneg, Potential.specificRelativeEntropy_eq_zero_of_mem_invariantG}
    \leanok

    For probability measures $\mu, \rho, \rho'$ and any volume $\Lambda$,
    $\mathcal H_\Lambda(\mu \mid \rho\gamma^\Phi_\Lambda) \leq \mathcal H_\Lambda(\mu \mid \rho'\gamma^\Phi_\Lambda) + 4\, t(\Lambda)$
    with $t(\Lambda)$ the boundary sum of (15.25). For $\mu \in \Prob_\Theta$ and any
    $\rho \in \mathcal G(\gamma^\Phi)$,
    $|\Lambda_n|^{-1}\mathcal H_{\Lambda_n}(\mu \mid \rho) \to k(\mu \mid \Phi) := P(\Phi) + \mu(f_\Phi) - h(\mu) \geq 0$
    (15.30)(b), (15.32), (15.35); for $\mu \in \mathcal G(\gamma^\Phi)$, without shift invariance,
    $|\Lambda_n|^{-1}\mathcal H_{\Lambda_n}(\mu \mid \mu\gamma^\Phi_{\Lambda_n}) \to 0$ (15.33); and
    $k(\mu \mid \Phi) = 0$ for $\mu \in \mathcal G_\Theta(\gamma^\Phi)$.
\end{proposition}
\begin{proof}
    \leanok
    (15.29) writes $\rho\gamma_\Lambda$ on $\mathcal F_\Lambda$ as a density against $\lambda^S$
    and (15.34) expresses $\mathcal H_\Lambda(\mu \mid \rho\gamma_\Lambda)$ through
    $\mathcal H_\Lambda(\mu)$, $\mu(H_\Lambda)$ and $\log Z_\Lambda$; the boundary estimate
    (15.25) and Theorems \ref{thm:sent-15-12}, \ref{thm:pr-15-30} give the limits.
\end{proof}

\begin{theorem}[Georgii (15.37)]
    \label{thm:vp-15-37}
    \uses{def:ql-quasilocal-spec, def:erg-invariant, def:gibbs-meas}
    \lean{MeasureTheory.GibbsMeasure.exists_subset_relativeEntropyIn_sub_le, MeasureTheory.GibbsMeasure.exists_density_of_relativeEntropyIn_sub_le, MeasureTheory.GibbsMeasure.isGibbsMeasure_of_forall_exists_density, MeasureTheory.GibbsMeasure.isGibbsMeasure_of_tendsto_relativeEntropyIn_div_card}
    \leanok

    Let $\gamma$ be quasilocal, $\mu, \nu$ shift-invariant probability measures with
    $\nu \in \mathcal G(\gamma)$, and suppose $|\Lambda_n|^{-1}\mathcal H_{\Lambda_n}(\mu \mid \nu) \to 0$
    along every family of boxes with all side lengths tending to infinity. Then
    $\mu \in \mathcal G(\gamma)$.
\end{theorem}
\begin{proof}
    \leanok
    Tile $[0, N(2R+1))^d$ by $N^d$ translates of $[-R,R]^d$; by monotonicity and shift
    invariance some translate $\Delta$ satisfies
    $\mathcal H_{\Delta \cup \Lambda}(\mu\mid\nu) - \mathcal H_{\Delta}(\mu\mid\nu) \leq \varepsilon$
    for the given $\Lambda$. The chain rule (15.38) turns this into a density
    $\tilde g$ of $\mu$ against $\nu$ on $\mathcal F_{\Delta\cup\Lambda}$ close to $1$ in
    $L^1$ relative to $\mathcal F_\Delta$; quasilocality of $\gamma$ and the DLR equation for
    $\nu$ then bound $|\mu(f) - \mu(\gamma_\Lambda f)|$ for local $f$ by $\varepsilon$-terms.
\end{proof}

\begin{theorem}[Variational principle, Georgii (15.39)]
    \label{thm:vp-15-39}
    \uses{prop:vp-15-30b, thm:vp-15-37}
    \lean{Potential.specificRelativeEntropy_eq_zero_iff_mem_invariantG, Potential.mem_invariantG_of_specificRelativeEntropy_eq_zero, Potential.specificEntropy_le_specificEnergy_add_pressure}
    \leanok

    For $\mu \in \Prob_\Theta$, $h(\mu) \leq \mu(f_\Phi) + P(\Phi)$, and, provided
    $\mathcal G_\Theta(\gamma^\Phi) \neq \emptyset$, equality holds iff
    $\mu \in \mathcal G_\Theta(\gamma^\Phi)$.
\end{theorem}
\begin{proof}
    \leanok
    Proposition \ref{prop:vp-15-30b} gives $k(\mu\mid\Phi) \geq 0$ with equality on
    $\mathcal G_\Theta(\gamma^\Phi)$; conversely $k(\mu\mid\Phi) = 0$ and
    $\rho \in \mathcal G_\Theta(\gamma^\Phi)$ give $|\Lambda_n|^{-1}\mathcal H_{\Lambda_n}(\mu\mid\rho) \to 0$
    by (15.30)(b), and Theorem \ref{thm:vp-15-37} applies.
\end{proof}
